\section{Algoritma Pembanding}

Sebagaimana ditetapkan pada variabel bebas jenis algoritma, CoDA-EDA dibandingkan dengan enam \textit{baseline}, yaitu \textit{compact-GA}, \textit{compact-PSO}, BOA, iBOA, ACO, dan UBK. 
Seluruh algoritma diimplementasikan dan dikonfigurasi dengan prinsip perbandingan yang setara agar hasil \textit{Evaluation} tidak dipengaruhi perbedaan representasi, anggaran komputasi, maupun kualitas implementasi.

\textit{Compact-GA} dan \textit{compact-PSO} mengikuti \textit{template} algoritma \textit{compact} dengan PV univariat, kompetisi pemenang–pecundang, dan ukuran populasi virtual (NP) sebagai parameter utama. 
Keduanya menggunakan representasi biner yang sama dengan CoDA-EDA agar perbedaan hasil terutama mencerminkan mekanisme pencarian dan pemodelan dependensi. 
\textit{Compact-PSO}, yang berasal dari ruang kontinu, diadaptasi ke representasi biner menggunakan pendekatan probabilistik yang sejalan dengan praktik pembineran metaheuristik kontinu \parencite{pan_survey_2023}.

BOA dan iBOA diimplementasikan menggunakan pembelajaran struktur berbasis BIC dan estimasi parameter kemungkinan maksimum tanpa pembatasan orde dependensi. 
BOA mempertahankan populasi eksplisit dan mempelajari ulang struktur pada setiap generasi, sedangkan iBOA memperbarui struktur dan parameter secara inkremental melalui skema turnamen. 
Karena karakteristik tersebut, iBOA menjadi pembanding terdekat untuk mengevaluasi mekanisme pembaruan CoDA-EDA pada H3.

ACO dan UBK digunakan sebagai \textit{baseline} metaheuristik berbasis populasi. 
ACO menggunakan mekanisme feromon dengan parameter evaporasi serta bobot heuristik $(\alpha,\beta)$, sedangkan UBK mengikuti mekanisme pencarian berbasis populasi sebagaimana publikasi aslinya. 
Parameter awal mengikuti rekomendasi sumber asli dan dapat dikalibrasi pada dataset penelitian agar kinerja \textit{baseline} tetap representatif.

Kesetaraan perbandingan dijaga melalui empat kontrol. 
Pertama, seluruh algoritma menggunakan representasi solusi dan fungsi \textit{fitness} yang sama. 
Kedua, anggaran komputasi disetarakan berdasarkan jumlah maksimum evaluasi \textit{fitness}, bukan jumlah iterasi, karena definisi iterasi berbeda antaralgoritma \parencite{rajwar_exhaustive_2023,omran_permutation_2022}. 
Anggaran ditetapkan sebagai fungsi dimensi $D$, mengingat dua dataset utama memiliki dimensi berbeda, yaitu 436 pada \textit{workforce} dan 1.907 pada \textit{e-document}. 
Ketiga, anggaran dihitung untuk keseluruhan proses \textit{separate-and-conquer}, bukan per aturan. 
Keempat, seluruh algoritma dijalankan pada perangkat keras dan dataset yang sama untuk setiap skenario, dengan replikasi menggunakan \textit{random seed} berbeda untuk mengakomodasi variabilitas stokastik.

\begin{longtable}{|c|>{\raggedright\arraybackslash}p{1.7cm}|>{\raggedright\arraybackslash}p{2.2cm}|>{\raggedright\arraybackslash}p{2cm}|>{\raggedright\arraybackslash}p{2cm}|>{\raggedright\arraybackslash}p{3.2cm}|}
\caption{Algoritma Pembanding dan Perlakuan Kesetaraan Implementasi}
\label{tab:algoritma-pembanding} \\
\hline
\textbf{No} & \textbf{Algoritma} & \textbf{Sumber} & \textbf{Kelas} & \textbf{Parameter Kunci} & \textbf{Catatan Kesetaraan} \\
\hline
\endfirsthead

\multicolumn{6}{c}{{\bfseries \tablename\ \thetable{} -- lanjutan dari halaman sebelumnya}} \\
\hline
\textbf{No} & \textbf{Algoritma} & \textbf{Sumber} & \textbf{Kelas} & \textbf{Parameter Kunci} & \textbf{Catatan Kesetaraan} \\
\hline
\endhead

\hline
\multicolumn{6}{r}{Lanjut ke halaman berikutnya} \\
\endfoot

\hline
\endlastfoot

% ==================== DATA BARIS ====================
1 & \textit{compact-GA} & \parencite{khalfi_metaheuristics_2023} & \textit{Compact}, univariat (diskret) & Ukuran populasi virtual (NP) & Representasi biner identik dengan CoDA-EDA; tanpa pemodelan dependensi \\
\hline
2 & \textit{compact-PSO} & \parencite{khalfi_metaheuristics_2023} & \textit{Compact}, univariat (kontinu, diadaptasi) & NP; parameter kecepatan-posisi & Representasi diadaptasi ke ruang diskret agar sebanding dengan \textit{compact-GA} \\
\hline
3 & BOA & \parencite{pelikan_iboa_2008} & \textit{Population-based}, dependensi tak terbatas & Ukuran populasi; ambang BIC & Struktur dipelajari ulang penuh tiap generasi; representasi biner sama \\
\hline
4 & iBOA & \parencite{pelikan_iboa_2008} & Inkremental, dependensi tak terbatas & Ambang/ kriteria koneksi struktur & Skema turnamen sama dengan CoDA-EDA; pembanding langsung untuk H3 \\
\hline
5 & ACO & \parencite{shang_abac_2024} & \textit{Population-based}, \textit{swarm} & Jumlah semut; evaporasi feromon; $\alpha$\label{first:symbol-alpha}, $\beta$\label{first:symbol-beta} & Parameter mengikuti publikasi asli, dikalibrasi bila perlu \\
\hline
6 & UBK & \parencite{li_attribute_2026} & \textit{Population-based, swarm} & Ukuran populasi; parameter strategi berburu & Parameter mengikuti publikasi asli, dikalibrasi bila perlu \\
\hline

\end{longtable}

Selain kontrol eksperimen, kesetaraan juga dijaga pada operator bantu yang digunakan lintas algoritma. 
Operator \textit{repair} untuk memastikan aturan valid, yaitu memiliki sekurang-kurangnya satu aksi dan satu kondisi atribut aktif, serta prosedur pembenihan dari contoh \textit{permit} diterapkan identik pada seluruh tujuh kondisi algoritma. 
Seluruh solusi juga menggunakan tipe struktur data yang sama agar perbedaan konsumsi memori merefleksikan karakteristik algoritma, bukan \textit{overhead} implementasi. 
Kontrol ini penting terutama untuk menjaga kesahihan pengujian H2.

Sebelum eksperimen utama, setiap implementasi menjalani uji monotonisitas terhadap anggaran evaluasi. 
Peningkatan anggaran tidak seharusnya menghasilkan penurunan kualitas kebijakan secara sistematis. 
Uji ini digunakan untuk mendeteksi ketidakselarasan antara fungsi \textit{fitness} dan metrik evaluasi yang mungkin tidak terlihat pada satu tingkat anggaran. 
Implementasi yang menunjukkan pola penurunan konsisten akan diperiksa dan diperbaiki sebelum digunakan dalam pengujian hipotesis.